% ============================================================
%  FICHEIRO 1 — Modelação da Categoria FinSet
%
%  Conteúdo:
%    1. FinSet Clássica vs. Computacional
%    2. Objetos e Invariante de Representação
%    3. Morfismos (Identidade, Composição)
%    4. Funções Hash como Morfismos de FinSet
%    5. Homologia de Endomapas
%    6. Hash vs. Homologia: Análise Comparativa
%    7. Estrutura de Magma a partir de Endomapas
%    8. Classes de Conjugação em End(A)
% ============================================================

% ─────────────────────────────────────────────────────────────
\section{FinSet: Categoria Clássica vs.\ Computacional}
% ─────────────────────────────────────────────────────────────

A categoria \textbf{FinSet clássica} é fundamental na teoria das
categorias: os seus objetos são conjuntos finitos arbitrários e
os seus morfismos são funções totais entre eles. Nesta perspetiva
matemática abstrata, dois conjuntos com o mesmo número de elementos
são isomórficos, e a identidade de um conjunto não depende de
qualquer representação concreta dos seus elementos.

Ao contrário da categoria \textbf{Fin} (Secção~2.1), onde os objetos
são representados apenas pela sua cardinalidade (um natural $n$
representa o conjunto $\{1,\ldots,n\}$), a \textbf{FinSet Computacional}
representa cada objeto como uma \emph{matriz de MATLAB/Octave com
	linhas únicas e ordenadas lexicograficamente}. Esta representação
permite modelar conjuntos cujos elementos são registos multidimensionais
--- \textit{strings}, vetores inteiros, coordenadas no espaço ---
de forma direta e sem perda de informação estrutural.

\begin{table}[htbp]
	\centering
	\caption{Comparação entre FinSet Clássica e FinSet Computacional.}
	\begin{tabular}{lll}
		\toprule
		\textbf{Aspeto} & \textbf{FinSet Clássica} & \textbf{FinSet Computacional} \\
		\midrule
		Objetos        & Conjuntos finitos abstractos       & Matrizes MATLAB (linhas únicas, ordenadas) \\
		Morfismos      & Funções totais                     & Vetores de índices inteiros \\
		Identidade     & Abstracta (sem representação)      & Invariante de representação explícito \\
		Elementos      & Arbitrários                        & Linhas de matrizes numéricas ou char \\
		Composição     & $g \circ f$ (abstracta)            & \texttt{h = g(f)} (indexação) \\
		Isomorfismo    & Bijeção qualquer                   & Permutação de linhas \\
		\bottomrule
	\end{tabular}
	\label{tab:finset_classica_vs_comp}
\end{table}

A FinSet Computacional é o ambiente natural para implementar
construções universais (produto, coproduto, equalizador, pullback
e pushout), funções hash, homologia de endomapas e estruturas
algébricas como magmas, sendo todas estas noções rigorosamente
definíveis como morfismos ou funtores em FinSet.

% ─────────────────────────────────────────────────────────────
\section{Modelação Computacional da Categoria FinSet}
% ─────────────────────────────────────────────────────────────

\subsection{Objetos e Invariante de Representação}

Em MATLAB, um objeto da categoria FinSet computacional é
representado como uma matriz. Para garantir a definição
matemática de um conjunto (elementos únicos e sem ordem
intrínseca relevante para a igualdade), define-se o seguinte
invariante de representação.

\begin{definition}[Objeto de FinSet Computacional]
	Um \textbf{objeto} da categoria FinSet é uma variável
	MATLAB/Octave \texttt{A} que satisfaz o
	\emph{invariante de representação}:
	\begin{equation}
		\texttt{isequal(A,\;unique(A,\textquotesingle rows\textquotesingle))}
		\;=\; \texttt{true}.
		\label{eq:finset_inv}
	\end{equation}
	Isto significa que as linhas de \texttt{A} são distintas e estão
	ordenadas lexicograficamente. Cada linha representa um elemento
	do conjunto finito correspondente.
\end{definition}

\begin{remark}
	A função \texttt{unique(A,'rows')} do MATLAB realiza
	simultaneamente a remoção de duplicados e a ordenação
	lexicográfica. Qualquer operação que produza um novo objeto
	deve terminar com \texttt{sortrows(unique(X,'rows'))} para
	restaurar o invariante~\eqref{eq:finset_inv}.
\end{remark}

Esta abordagem permite modelar diversos tipos de conjuntos:
\begin{itemize}
	\item \textbf{Conjuntos escalares}: Vetores coluna de números
	inteiros.
	\item \textbf{Conjuntos de strings}: Matrizes de caracteres
	(\textit{char arrays}).
	\item \textbf{Produtos cartesianos}: Matrizes onde cada coluna
	representa uma dimensão do produto (e.g., pares de
	coordenadas).
	\item \textbf{Vetores multidimensionais}: Matrizes com colunas
	representando componentes (e.g., vetores em $\mathbb{Z}^3$).
\end{itemize}

\begin{lstlisting}[style=matlab, caption={Definição do invariante e exemplos de objetos em FinSet}]
	% Um objeto de FinSet e uma matriz MATLAB A tal que:
	%   isequal(A, unique(A,'rows')) == true
	is_FinSet_object = @(X) isequal(X, sortrows(unique(X,'rows')));
	
	% Objeto 1: conjunto de inteiros {1,...,7}
	A_int = (1:7)';
	assert(is_FinSet_object(A_int));
	
	% Objeto 2: conjunto de strings (alunos da sala)
	A_str = sortrows(['Abe'; 'Bar'; 'Car'; 'Dye'; 'Eve'; 'Fly';
	'Gym'; 'Hym'; 'Ive'; 'Jun'; 'Kay'; 'Lie';
	'May'; 'Neo'; 'Oli'; 'Pie'; 'Que'; 'Ray';
	'Sym'; 'Til']);
	assert(is_FinSet_object(A_str));
	
	% Objeto 3: mesas como pares (linha, coluna) em {1..5}x{1..3}
	[i, j] = ind2sub([5 3], 1:15);
	B_pairs = sortrows([i(:), j(:)]);   % 15 mesas
	assert(is_FinSet_object(B_pairs));
	
	% Objeto 4: vetores em Z^3 (inclui negativos)
	C_vec = sortrows([1 0 0; 0 1 0; 0 0 1; 0 0 0;
	-1 0 0; 0 -1 0; 0 0 -1]);
	assert(is_FinSet_object(C_vec));
	
	fprintf('A_int  : %d elementos\n', size(A_int,1));
	fprintf('A_str  : %d elementos\n', size(A_str,1));
	fprintf('B_pairs: %d elementos\n', size(B_pairs,1));
	fprintf('C_vec  : %d elementos\n', size(C_vec,1));
\end{lstlisting}

% ─────────────────────────────────────────────────────────────
\subsection{Morfismos}

\begin{definition}[Morfismo de FinSet]
	Um \textbf{morfismo} em FinSet é um triplo $(B,\, f,\, A)$ onde:
	\begin{itemize}
		\item $A$ e $B$ são objetos de FinSet;
		\item $f$ é um vetor de inteiros de comprimento $|A|$
		(i.e., \texttt{size(A,1)}) satisfazendo:
		\begin{equation}
			\texttt{numel(f) == size(A,1)}
			\quad\text{e}\quad
			f(i) \in \{1, \ldots, |B|\},\;\forall i.
			\label{eq:finset_morph}
		\end{equation}
	\end{itemize}
	A condição \texttt{numel(f) == size(A,1)} garante que $f$ é uma
	\emph{função total}: todo o elemento do domínio tem exatamente
	uma imagem. A expressão \texttt{B(f,:)} devolve, para cada linha
	$i$ de $A$, a linha $f(i)$ de $B$, codificando assim $f: A \to B$.
\end{definition}

\begin{lstlisting}[style=matlab, caption={Validação de um morfismo em FinSet}]
	function ok = is_FinSet_morphism(B, f, A)
	% (B, f, A): morfismo f: A -> B em FinSet
	% f e vetor coluna de indices com:
	%   nnz(f) == size(A,1)   [total: sem zeros]
	%   f(i) in {1..size(B,1)}
	f = f(:);
	ok = (numel(f)   == size(A,1)) && ...
	(nnz(f)     == size(A,1)) && ...
	all(f >= 1)               && ...
	all(f <= size(B,1));
	end
	
	% Morfismo: atribuicao de mesas (20 alunos -> 15 mesas, ciclico)
	f_mesas = mod((1:20)'-1, 13) + 1;
	assert(is_FinSet_morphism(B_pairs, f_mesas, A_str));
	
	% A expressao B(f,:) devolve a mesa de cada aluno
	mesas_alunos = B_pairs(f_mesas, :);
	fprintf('Primeiros 5 alunos e as suas mesas:\n');
	for k = 1:5
	fprintf('  %s -> mesa (%d,%d)\n', A_str(k,:), ...
	mesas_alunos(k,1), mesas_alunos(k,2));
	end
	
	% Morfismo identidade
	id_A = (1:size(A_int,1))';
	assert(is_FinSet_morphism(A_int, id_A, A_int));
\end{lstlisting}

\begin{definition}[Identidade e Composição]
	A \textbf{identidade} em $A$ é o morfismo $(A,\;(1:n)',\;A)$
	com $n = |A|$.
	A \textbf{composição} de $(C,g,B)$ com $(B,f,A)$ é
	$(C,\; g(f),\; A)$, i.e., a indexação composta \texttt{h = g(f)}.
\end{definition}

\begin{proposition}
	A composição em FinSet é associativa e a identidade é neutra,
	fazendo de FinSet uma categoria.
\end{proposition}
\begin{proof}
	Para $(D,h,C)$, $(C,g,B)$, $(B,f,A)$:
	$(h \circ g) \circ f = h(g(f)) = h \circ (g \circ f)$,
	pois a indexação de vetores é associativa em MATLAB.
\end{proof}

\begin{lstlisting}[style=matlab, caption={Composição de morfismos e verificação de associatividade}]
	function h = compose_FinSet(g, f)
	% Composicao: (C,g,B) o (B,f,A) = (C, g(f), A)
	% h(i) = g(f(i))
	h = g(f);
	end
	
	% Verificacao de associatividade: (h o g) o f == h o (g o f)
	f1 = [2;1;3;2];   % A(4) -> B(3)
	g1 = [3;1;2];     % B(3) -> C(3)
	h1 = [2;1;3];     % C(3) -> D(3)
	assert(isequal(compose_FinSet(h1, compose_FinSet(g1,f1)), ...
	compose_FinSet(compose_FinSet(h1,g1), f1)));
	fprintf('Composicao associativa: OK\n');
\end{lstlisting}

% ─────────────────────────────────────────────────────────────
\subsection{Funções Hash como Morfismos de FinSet}

\begin{definition}[Função Hash em FinSet]
	Uma \textbf{função hash} $h: U \to \{1,\ldots,m\}$ é um morfismo
	em FinSet de um universo de chaves $U$ (potencialmente grande)
	para um conjunto de \emph{buckets} de cardinalidade $m$ pequena.
	Em notação de FinSet, $h = (B_m,\, f_h,\, U)$ com
	$B_m = (1:m)'$.
\end{definition}

A diferença fundamental entre uma função hash e um morfismo
injetivo é a presença de \textbf{colisões}: $h(x) = h(y)$
com $x \neq y$. As colisões ocorrem inevitavelmente, sendo uma
consequência direta da perda de injetividade (Teorema de Dirichlet).
A gestão de colisões é uma escolha adicional (encadeamento,
endereçamento aberto, etc.) que corresponde categoricamente à
escolha de uma \emph{secção} do morfismo de quociente associado.

\begin{definition}[Colisão e Resolução]
	Dada $h: U \to B_m$, diz-se que $x, y \in U$ com $x \neq y$
	\textbf{colidem} se $h(x) = h(y)$.
	A \textbf{resolução por encadeamento} associa a cada bucket
	$b \in B_m$ uma lista $T[b] = \{x \in U : h(x) = b\}$,
	que corresponde à \emph{fibra} de $h$ sobre $b$.
\end{definition}

\begin{lstlisting}[style=matlab, caption={Três famílias de funções hash como morfismos de FinSet}]
	%% 1) Hash por divisao: h(k) = mod(k,m) + 1
	%%    Morfismo: U = Z -> {1,...,m}
	function idx = hash_division(key, m)
	idx = mod(key, m) + 1;
	end
	
	%% 2) Hash por multiplicacao (Knuth):
	%%    A = (sqrt(5)-1)/2 (razao aurea conjugada)
	%%    h(k) = floor(m * frac(k*A))
	function idx = hash_multiplication(key, m)
	A_knuth = (sqrt(5) - 1) / 2;
	idx = floor(m * mod(key * A_knuth, 1)) + 1;
	end
	
	%% 3) Hash polinomial para strings:
	%%    h(s) = (sum s[i]*p^i) mod m
	function idx = hash_polynomial(key_str, m)
	p = 31; val = 0;
	for i = 1:numel(key_str)
	val = mod(val * p + double(key_str(i)), m);
	end
	idx = val + 1;
	end
	
	%% Demonstracao: U = {1,...,20}, m = 7 buckets
	U = 1:20; m = 7;
	h_div  = arrayfun(@(k) hash_division(k, m),       U);
	h_mult = arrayfun(@(k) hash_multiplication(k, m),  U);
	
	fprintf('Chaves : '); fprintf('%3d', U);     fprintf('\n');
	fprintf('Div    : '); fprintf('%3d', h_div);  fprintf('\n');
	fprintf('Mult   : '); fprintf('%3d', h_mult); fprintf('\n');
	
	% Contagem de colisoes por bucket
	for b = 1:m
	fprintf('Bucket %d: Div=%d | Mult=%d colisoes\n', ...
	b, sum(h_div==b), sum(h_mult==b));
	end
\end{lstlisting}

\begin{lstlisting}[style=matlab, caption={Tabela hash com encadeamento e aplicação a ficheiros}]
	%% Tabela Hash com Encadeamento (Chaining)
	function T = build_hash_table(keys, m, hash_fn)
	T = cell(m, 1);
	for i = 1:numel(keys)
	b = hash_fn(keys(i));
	T{b} = [T{b}, keys(i)];
	end
	end
	
	T_div = build_hash_table(U, m, @(k) hash_division(k,m));
	fprintf('Tabela Hash (divisao, m=%d):\n', m);
	for b = 1:m
	fprintf('  Bucket %d: [%s]\n', b, num2str(T_div{b}));
	end
	
	%% Hash de conteudo para ficheiros (simula hash de bytes)
	function digest = hash_file_content(data, m)
	p = 257; digest = 0;
	for i = 1:numel(data)
	digest = mod(digest * p + data(i), m);
	end
	digest = digest + 1;
	end
	
	%% Simulacao: 5 "ficheiros" como vetores de bytes
	files = {[72 101 108 108 111],    ...  % "Hello"
		[87 111 114 108 100],    ...  % "World"
		[77 65 84 76 65 66],     ...  % "MATLAB"
		[70 105 110 83 101 116], ...  % "FinSet"
		[72 101 108 108 111]};        % "Hello" (duplicado)
	m_files = 8;
	fprintf('Hash de ficheiros (m=%d buckets):\n', m_files);
	for i = 1:numel(files)
	h = hash_file_content(files{i}, m_files);
	fprintf('  Ficheiro %d -> bucket %d\n', i, h);
	end
	fprintf('(Ficheiros 1 e 5 colidem: conteudo identico!)\n');
\end{lstlisting}

% ─────────────────────────────────────────────────────────────
\subsection{Homologia de Endomapas}

\begin{definition}[Homologia de um Endomapa]
	Dado $\varphi: A \to A$ com $|A| = n$, o
	\textbf{grafo dirigido associado} $G(\varphi)$ tem vértices
	$A$ e arestas $i \to \varphi(i)$.
	A \textbf{homologia} de $\varphi$ é a matriz
	$H(\varphi) \in \mathbb{N}^{c \times 4}$,
	onde $c$ é o número de componentes conexas (no sentido do
	grafo subjacente não dirigido) e cada linha~$r$ contém:
	\begin{equation}
		H(\varphi)_{r} = \bigl(\#\text{grau-in }0,\;
		\#\text{grau-in }1,\;
		\#\text{grau-in }\geq 2,\;
		\text{comp.\ do ciclo}\bigr)_r.
		\label{eq:homology}
	\end{equation}
\end{definition}

\begin{remark}
	Todo o endomapa $\varphi$ de um conjunto finito possui pelo
	menos um ciclo em cada componente conexa (Teorema de Rho de
	Floyd). Os vértices \emph{fora} do ciclo formam ``caudas''
	(grau de entrada 0 no ciclo). Os vértices de grau de entrada
	$\geq 2$ são pontos de ramificação.
\end{remark}

O código analisa a ``homologia'' do grafo dirigido, determinando
para cada componente conexa:
\begin{enumerate}
	\item Número de elementos terminais (\textit{tails}) ---
	vértices com grau de entrada 0.
	\item Vértices de ramificação (grau de entrada $\geq 2$).
	\item O comprimento invariável do ciclo central.
\end{enumerate}

\begin{lstlisting}[style=matlab, caption={Cálculo da matriz de homologia de um endomapa}]
	function H = endomap_homology(phi)
	% phi: vetor representando phi: {1..n} -> {1..n}
	% Devolve a matriz de homologia H (uma linha por componente)
	phi = phi(:)'; n = numel(phi);
	
	% --- Passo 1: Componentes conexas (Union-Find com path compression)
	parent = 1:n;
	for i = 1:n
	ra = i;
	while parent(ra) ~= ra
	parent(ra) = parent(parent(ra)); ra = parent(ra);
	end
	rb = phi(i);
	while parent(rb) ~= rb
	parent(rb) = parent(parent(rb)); rb = parent(rb);
	end
	if ra ~= rb, parent(ra) = rb; end
	end
	comp_label = zeros(1, n);
	for i = 1:n
	r = i;
	while parent(r) ~= r, r = parent(r); end
	comp_label(i) = r;
	end
	[comp_ids, ~, comp_idx] = unique(comp_label);
	comp_idx = comp_idx(:)';   % forcar vetor linha
	n_comp = numel(comp_ids);
	
	% --- Passo 2: Grau de entrada de cada vertice
	in_deg = histc(phi, 1:n);  % in_deg(i) = #{j : phi(j) = i}
	
	% --- Passo 3: Detetar ciclos por visita iterada
	on_cycle = false(1, n);
	for start = 1:n
	visited = false(1, n); x = start;
	while ~visited(x) && ~on_cycle(x)
	visited(x) = true; x = phi(x);
	end
	if on_cycle(x), continue; end
	if visited(x)
	y = x;
	while ~on_cycle(y), on_cycle(y) = true; y = phi(y); end
	end
	end
	
	% --- Passo 4: Montar matriz de homologia
	H = zeros(n_comp, 4);
	for c = 1:n_comp
	mask_c = (comp_idx == c);
	ideg_c = in_deg(mask_c);
	H(c,1) = sum(ideg_c == 0);           % grau 0 (caudas)
	H(c,2) = sum(ideg_c == 1);           % grau 1
	H(c,3) = sum(ideg_c >= 2);           % grau >= 2 (ramif.)
	H(c,4) = sum(mask_c & on_cycle);     % comprimento ciclo
	end
	H = sortrows(H);   % forma canonica
	end
\end{lstlisting}

\begin{lstlisting}[style=matlab, caption={Exemplos de homologia do quadro da aula}]
	%% Exemplo 1 (quadro, Imagem 4):
	%% phi: 1->2, 2->1, 3->4, 4->6, 5->6, 6->3, 7->7
	phi_ex1 = [2, 1, 4, 6, 6, 3, 7];
	H1 = endomap_homology(phi_ex1);
	fprintf('Homologia phi1 (3 componentes esperadas):\n');
	fprintf('  grau0 | grau1 | grau>=2 | ciclo\n');
	for r = 1:size(H1,1)
	fprintf('    %2d  |  %2d   |   %2d    |  %2d\n', ...
	H1(r,1), H1(r,2), H1(r,3), H1(r,4));
	end
	%% Resultado esperado:
	%%   grau0 grau1 grau2 ciclo
	%%     0     2     0     2    (ciclo {1,2})
	%%     1     2     1     3    (ciclo {3,4,6}, cauda {5})
	%%     0     1     0     1    (ciclo {7})
	
	%% Exemplo 2 (quadro, Imagem 2, seccao esquerda):
	%% phi: 1->5, 2->3, 3->2, 4->5, 5->6, 6->7, 7->4
	phi_ex2 = [5, 3, 2, 5, 6, 7, 4];
	H2 = endomap_homology(phi_ex2);
	fprintf('Homologia phi2:\n'); disp(H2);
	
	%% Exemplo 3 (quadro, Imagem 3, seccao direita):
	%% phi: 1->2, 2->3, 3->4, 4->5, 5->2, 6->7, 7->7
	phi_ex3 = [2, 3, 4, 5, 2, 7, 7];
	H3 = endomap_homology(phi_ex3);
	fprintf('Homologia phi3:\n'); disp(H3);
\end{lstlisting}

\begin{proposition}
	A homologia $H: \mathrm{End}(A) \to \mathrm{FinSet}$ é um
	\textbf{invariante de conjugação}: se $\varphi \sim \psi$
	(i.e., existe bijeção $w : A \to A$ tal que
	$\varphi \circ w = w \circ \psi$), então
	$H(\varphi) = H(\psi)$.
\end{proposition}
\begin{proof}
	A conjugação por $w$ é um isomorfismo de grafos de $G(\psi)$
	em $G(\varphi)$, preservando assim a estrutura das componentes
	conexas, os graus de entrada e os comprimentos de ciclo.
\end{proof}

% ─────────────────────────────────────────────────────────────
\subsection{Hash vs.\ Homologia: Análise Comparativa}

Ambas as construções são morfismos de FinSet que \emph{resumem}
informação de um objeto grande (universo $U$ ou $\mathrm{End}(A)$)
num objeto mais pequeno. A diferença essencial está no objetivo:
eficiência de acesso (hash) vs.\ invariância estrutural
(homologia).

\begin{table}[htbp]
	\centering
	\caption{Comparação entre funções hash e a função de homologia.}
	\begin{tabular}{lll}
		\toprule
		\textbf{Propriedade} & \textbf{Função Hash} & \textbf{Homologia} \\
		\midrule
		Domínio        & $U$ (chaves arbitrárias)         & $\mathrm{End}(A)$ (endomapas) \\
		Codomínio      & $\{1,\ldots,m\}$ (buckets)       & Matrizes de $\mathbb{N}^{c\times 4}$ \\
		Objetivo       & Indexação rápida                  & Classificação estrutural \\
		Colisões       & Inevitáveis (pretendidas)         & Possíveis (indesejadas) \\
		Complexidade   & $O(1)$ (avaliação)               & $O(n)$ (cálculo) \\
		Invariância    & Nenhuma garantida                 & Sob conjugação \\
		Aplicação      & BD, ficheiros, criptografia       & Classes de $\mathrm{End}(A)$ \\
		\bottomrule
	\end{tabular}
	\label{tab:hash_vs_hom}
\end{table}

Em particular, a homologia pode ser usada como uma
\textbf{hash de classificação} para endomapas: ao invés de
comparar dois endomapas diretamente (custo $O(n)$), calcula-se
$H(\varphi)$ e $H(\psi)$ e comparam-se as matrizes resultantes.
Se $H(\varphi) \neq H(\psi)$, os endomapas não são conjugados
(e.g., num ficheiro de base de dados de estruturas, basta
indexar pelo hash da homologia).

\begin{lstlisting}[style=matlab, caption={Uso da homologia como hash de classificação de endomapas}]
	function same = same_homology(phi1, phi2)
	H1 = endomap_homology(phi1);
	H2 = endomap_homology(phi2);
	same = isequal(sortrows(H1), sortrows(H2));
	end
	
	%% Indexacao de endomapas por homologia
	phi_list = {[2,1,3,4,5], [1,2,3,4,5], [2,1,4,3,5], ...
		[3,1,2,4,5], [2,3,1,4,5], [2,1,5,3,4]};
	
	groups = containers.Map('KeyType','char','ValueType','any');
	for k = 1:numel(phi_list)
	H   = endomap_homology(phi_list{k});
	key = mat2str(H);     % string da matriz = chave hash
	if groups.isKey(key)
	groups(key) = [groups(key), k];
	else
	groups(key) = k;
	end
	end
	disp('Grupos de conjugacao por homologia:');
	keys_g = groups.keys;
	for g = 1:numel(keys_g)
	fprintf('  Homologia %s -> endomapas %s\n', ...
	keys_g{g}, mat2str(groups(keys_g{g})));
	end
\end{lstlisting}

% ─────────────────────────────────────────────────────────────
\subsection{Estrutura de Magma a partir de Endomapas}

\begin{definition}[Magma derivado de um Endomapa]
	Dado $\varphi: A \to A$, define-se a operação binária
	$*_\varphi$ em $A$ por:
	\begin{equation}
		x *_\varphi x = \varphi(x)
		\qquad\text{e}\qquad
		x^{n+1} = \varphi^n(x),
		\label{eq:magma_endo}
	\end{equation}
	onde $\varphi^n$ denota a $n$-ésima iteração de $\varphi$.
	Para $x \neq y$, a regra geral é:
	$x *_\varphi y = x$ ou $y$ dependendo do tamanho da componente
	(conforme o quadro da aula: elementos de componentes maiores
	absorvem os menores).
	A tabela de Cayley $M_\varphi$ de $(A, *_\varphi)$ é uma
	matriz $|A| \times |A|$ de inteiros.
\end{definition}

\begin{lstlisting}[style=matlab, caption={Construção da tabela de Cayley do magma de um endomapa}]
	function [M_cayley, phi_iter] = endomap_to_magma(phi)
	% Constroi a tabela de Cayley do magma derivado de phi
	% Regra: x*x = phi(x); x*y = x se i<j, senao j
	phi = phi(:)'; n = numel(phi);
	
	% Calcular iteradas: phi_iter{k}(i) = phi^k(i)
	max_iter = 2*n;
	phi_iter = cell(1, max_iter+1);
	phi_iter{1} = 1:n;     % phi^0 = identidade
	for k = 1:max_iter
	phi_iter{k+1} = phi(phi_iter{k});
	end
	
	% Tabela de Cayley
	M_cayley = zeros(n, n);
	for i = 1:n
	for j = 1:n
	if i == j
	M_cayley(i,j) = phi(i);  % x*x = phi(x)
	elseif i < j
	M_cayley(i,j) = i;       % absorcao pelo menor
	else
	M_cayley(i,j) = j;
	end
	end
	end
	end
	
	%% Exemplo do quadro (Imagem 5): phi: 1->4, 2->3, 3->5, 4->1, 5->2
	phi_mag = [4, 3, 5, 1, 2];
	[M_mag, phi_iter_mag] = endomap_to_magma(phi_mag);
	fprintf('phi = [%s]\n', num2str(phi_mag));
	fprintf('Tabela de Cayley do magma:\n');
	fprintf('    ');
	for j = 1:size(M_mag,2), fprintf('%3d', j); end
	fprintf('\n');
	for i = 1:size(M_mag,1)
	fprintf('%3d ', i);
	fprintf('%3d', M_mag(i,:));
	fprintf('\n');
	end
	
	%% Verificacao: x^2 = x*x = phi(x)
	assert(all(diag(M_mag)' == phi_mag));
	fprintf('x*x = phi(x) para todo x: OK\n');
\end{lstlisting}

\begin{proposition}
	A construção $(A, \varphi) \mapsto (A, *_\varphi)$ é
	\textbf{functorial}: um morfismo $f: (A_1, \varphi_1) \to
	(A_2, \varphi_2)$ em \textbf{End}
	(i.e., $f \circ \varphi_1 = \varphi_2 \circ f$) induz um
	homomorfismo $f: (A_1, *_{\varphi_1}) \to (A_2, *_{\varphi_2})$
	em \textbf{Mag}, satisfazendo
	$f(x *_{\varphi_1} y) = f(x) *_{\varphi_2} f(y)$.
\end{proposition}

\begin{lstlisting}[style=matlab, caption={Verificação de functorialidade End -> Mag}]
	function ok = is_magma_morphism(M1, M2, f)
	% f: A1 -> A2 e morfismo de magmas se
	% f(x*y) = f(x)*f(y) para todo x,y
	n = size(M1,1); ok = true;
	for i = 1:n
	for j = 1:n
	lhs = f(M1(i,j));       % f(i*j)
	rhs = M2(f(i), f(j));   % f(i)*f(j)
	if lhs ~= rhs, ok = false; return; end
	end
	end
	end
	
	%% Morfismo de endomapas induz morfismo de magmas?
	phi1_m = [2,1,4,3,5]; phi2_m = [2,1,3,5,4];
	f_end  = [1,2,3,4,5]; % identidade (caso trivial)
	M1_m   = endomap_to_magma(phi1_m);
	M2_m   = endomap_to_magma(phi2_m);
	ok_comm = all(f_end(phi1_m) == phi2_m(f_end));
	ok_mag  = is_magma_morphism(M1_m, M2_m, f_end);
	fprintf('Morfismo de End -> Mag:\n');
	fprintf('  Comuta com phi: %d\n', ok_comm);
	fprintf('  E morfismo de magmas: %d\n', ok_mag);
\end{lstlisting}

% ─────────────────────────────────────────────────────────────
\subsection{Classes de Conjugação em End(A)}

\begin{definition}[Conjugação em End(A)]
	Dados $u, v \in \mathrm{End}(A)$, diz-se que $u$ e $v$ são
	\textbf{conjugados}, escrevendo $u \sim v$, se existe
	$w \in \mathrm{Bij}(A)$ tal que $u \circ w = w \circ v$.
\end{definition}

A relação $\sim$ é de equivalência. A homologia fornece um
invariante para separar classes: se $H(u) \neq H(v)$,
então $u \not\sim v$.

\begin{lstlisting}[style=matlab, caption={Algoritmo de verificação de conjugação em End(A)}]
	function ok = are_conjugate(phi1, phi2)
	% Verifica se phi1 ~ phi2
	% phi1 ~ phi2 <=> existe permutacao w: phi1(w(x)) = w(phi2(x))
	n = numel(phi1);
	% Condicao necessaria: mesma homologia
	if ~isequal(sortrows(endomap_homology(phi1)), ...
	sortrows(endomap_homology(phi2)))
	ok = false; return;
	end
	% Forca bruta para n pequeno: testar todas as permutacoes
	perms_w = perms(1:n);
	ok = false;
	for k = 1:size(perms_w,1)
	w = perms_w(k,:);
	if isequal(phi1(w), w(phi2))   % u o w = w o v
	ok = true; return;
	end
	end
	end
	
	%% Exemplos: A = {1,2,3,4,5}
	phi_a = [2, 1, 4, 3, 5];   % dois 2-ciclos + ponto fixo
	phi_b = [2, 1, 5, 3, 4];   % outro endomapa com 2 ciclos
	phi_c = [1, 2, 3, 4, 5];   % identidade
	
	fprintf('phi_a ~ phi_b? %d\n', are_conjugate(phi_a, phi_b));
	fprintf('phi_a ~ phi_c (id)? %d\n', are_conjugate(phi_a, phi_c));
	
	%% Agrupamento por homologia (invariante de conjugacao)
	phi_list = {[2,1,3,4,5], [1,2,3,4,5], [2,1,4,3,5], ...
		[3,1,2,4,5], [2,3,1,4,5], [2,1,5,3,4]};
	fprintf('\nGrupos por homologia (invariante de conjugacao):\n');
	homologies = {};
	for k = 1:numel(phi_list)
	H = endomap_homology(phi_list{k});
	h_str = mat2str(H);
	found = false;
	for g = 1:numel(homologies)
	if strcmp(homologies{g}.key, h_str)
	homologies{g}.members(end+1) = k;
	found = true; break;
	end
	end
	if ~found
	homologies{end+1} = struct('key', h_str, 'members', k);
	end
	end
	for g = 1:numel(homologies)
	fprintf('  Grupo %d: endomapas %s\n', g, ...
	mat2str(homologies{g}.members));
	end
\end{lstlisting}